\chapter{Physics-Informed Neural Networks}
\label{ch:pinn}

PINN tidak punya mata kuliahnya sendiri dalam perjalanan saya. Ia muncul sedikit-sedikit: sebagai
neural field yang diatur persamaan di akhir kuliah machine learning lanjut (\cref{ch:mlat}), sebagai
kemampuan autodiff menurunkan keluaran terhadap input di kuliah deep learning dasar (\cref{ch:dlbasic}),
dan sebagai salah satu pendekatan yang terus muncul dalam bacaan tentang machine learning untuk
mekanika fluida selama mengerjakan tugas akhir. Bab ini merangkai potongan-potongan itu.

Gagasan PINN sederhana. Alih-alih menyelesaikan persamaan diferensial di grid, kita mewakili solusinya
dengan jaringan saraf dan melatih jaringan itu agar memenuhi persamaannya. Gagasan ini sudah dicoba pada
akhir 1990-an \citep{lagaris1998}, dan dipopulerkan kembali dengan nama physics-informed neural networks
oleh \citet{raissi2019}, setelah autodiff dan GPU membuatnya praktis.

\section{Resep dasar}

Ambil persamaan panas satu dimensi, $\partial_tu=\kappa\,\partial_{xx}u$ pada $x\in[0,1]$ dan
$t\in[0,T]$, dengan syarat awal $u(x,0)=u_0(x)$ dan syarat batas $u(0,t)=u(1,t)=0$. Sebuah MLP
$u_\theta(x,t)$ menerima koordinat dan mengeluarkan temperatur. Untuk titik mana pun, autodiff memberi
$\partial_tu_\theta$ dan $\partial_{xx}u_\theta$ secara eksak, sehingga kita bisa menghitung
\istilah{residual} persamaan,
\begin{equation}
r_\theta(x,t)=\partial_tu_\theta(x,t)-\kappa\,\partial_{xx}u_\theta(x,t).
\end{equation}
Solusi sejati membuat residual nol di mana-mana. Maka kita meminimalkan jumlah kuadrat residual di
sekumpulan titik acak di dalam domain, yang disebut \istilah{collocation points}, ditambah suku untuk syarat
awal dan syarat batas:
\begin{equation}
\Loss(\theta)=\frac{\lambda_r}{N_r}\sum_{i=1}^{N_r}r_\theta(x_i,t_i)^2
+\frac{\lambda_0}{N_0}\sum_{j=1}^{N_0}\big(u_\theta(x_j,0)-u_0(x_j)\big)^2
+\frac{\lambda_b}{N_b}\sum_{k=1}^{N_b}u_\theta(x_k^b,t_k)^2.
\label{eq:pinnloss}
\end{equation}
Kalau ada pengukuran, suku keempat menambahkan selisih kuadrat antara $u_\theta$ dan data di titik-titik
pengukuran. Bobot $\lambda$ menentukan seberapa penting masing-masing suku, dan seperti akan kita lihat,
memilihnya ternyata salah satu bagian tersulit. \Cref{fig:pinn} merangkum alurnya.

\begin{figure}[ht]
\centering
\begin{tikzpicture}[node distance=0.5cm and 0.7cm]
  \node[kotak] (in) {$(x,t)$};
  \node[kotak=biru,right=of in] (net) {MLP $u_\theta$};
  \node[kotak,right=of net] (u) {$u_\theta(x,t)$};
  \node[kotak=ungu,above right=0.2cm and 0.7cm of u] (ad) {autodiff:\\$\partial_tu_\theta,\ \partial_{xx}u_\theta$};
  \node[kotak=hijau,below right=0.2cm and 0.7cm of u] (bc) {syarat awal,\\batas, data};
  \node[kotak=oranye,right=3.4cm of u] (L) {$\Loss(\theta)$};
  \draw[panah] (in) -- (net);
  \draw[panah] (net) -- (u);
  \draw[panah] (u) |- (ad);
  \draw[panah] (u) |- (bc);
  \draw[panah] (ad) -| node[label,above,pos=0.25]{residual} (L);
  \draw[panah] (bc) -| (L);
  \draw[panah,dashed] (L.south) -- ++(0,-1.2) -| node[label,below,pos=0.25]{gradien terhadap $\theta$} (net.south);
\end{tikzpicture}
\caption{Alur PINN. Jaringan menerima koordinat, autodiff memberi turunan yang dibutuhkan persamaan, dan loss
menggabungkan residual persamaan dengan syarat awal, syarat batas, dan data.}
\label{fig:pinn}
\end{figure}

Turunan terhadap input dihitung dengan autodiff yang sama dengan backpropagation (\cref{ch:dlbasic}), hanya
dirantai sekali atau dua kali lagi. Untuk satu neuron $u=\tanh(wx+b)$, turunan pertamanya
$u_x=w\,(1-u^2)$ dan turunan keduanya $u_{xx}=-2w^2u\,(1-u^2)$. Contoh kecil ini menjelaskan satu aturan praktis
PINN: fungsi aktivasi harus mulus. Dengan ReLU, turunan kedua bernilai nol hampir di mana-mana, sehingga
jaringan tidak bisa ``merasakan'' suku difusi sama sekali, dan residual persamaan panas hanya melihat
$\partial_tu$. Karena itu PINN hampir selalu memakai tanh, sinus, atau aktivasi mulus lain.

Mengapa ini masuk akal? Dari sudut pandang \cref{ch:data}, suku residual adalah prior: di antara semua
fungsi yang cocok dengan data dan syarat batas, pilih yang memenuhi hukum fisika. Dari sudut pandang
\cref{ch:mlat}, $u_\theta$ adalah neural field, dan persamaan diferensial adalah aturan yang ditegakkan
pada field itu, persis seperti persamaan Eikonal pada signed distance field.

\section{Contoh kecil yang bisa dihitung tangan}

Supaya terasa konkret, ambil persamaan peluruhan $\dot u=-ku$ dengan $u(0)=1$ dan $k=1$, pada
$t\in[0,1]$. Solusi sejatinya $e^{-t}$, sehingga $u(1)=e^{-1}\approx0{,}368$. \citet{lagaris1998}
menyarankan bentuk solusi yang memenuhi syarat awal secara otomatis, $u(t)=1+t\,N(t)$, dengan $N$ sebuah
jaringan. Untuk contoh tangan, kita ganti jaringan itu dengan konstanta $a$ saja, sehingga
$u(t)=1+at$. Residualnya
\begin{equation}
r(t)=\dot u+u=a+1+at=a(1+t)+1.
\end{equation}
Kita cari $a$ yang meminimalkan $J(a)=\int_0^1r(t)^2\dd t$, versi kontinu dari suku residual dalam
\cref{eq:pinnloss}. Turunan terhadap $a$ adalah $2\int_0^1\big(a(1+t)+1\big)(1+t)\dd t$. Menyamakannya
dengan nol memberi
\begin{equation}
a\int_0^1(1+t)^2\dd t=-\int_0^1(1+t)\dd t
\quad\Longrightarrow\quad a\cdot\frac73=-\frac32
\quad\Longrightarrow\quad a=-\frac9{14}\approx-0{,}643.
\end{equation}
Maka $u(1)\approx1-0{,}643=0{,}357$, hanya sekitar tiga persen di bawah nilai sebenarnya, padahal
``jaringan''-nya hanya satu bilangan. Jaringan sungguhan dengan ribuan parameter bisa mendekati
lengkungan eksponensial dengan jauh lebih baik. Contoh ini memperlihatkan inti PINN: tidak ada data
solusi sama sekali, hanya persamaan dan syarat awal, dan meminimalkan residual sudah cukup untuk
menemukan solusi yang masuk akal.

\section{Masalah invers}

Kekuatan PINN yang paling meyakinkan terletak pada masalah invers. Misalkan konduktivitas $\kappa$ tidak
diketahui, tetapi kita punya beberapa pengukuran temperatur. Cukup jadikan $\kappa$ parameter yang ikut
dilatih bersama bobot jaringan. Suku data memaksa $u_\theta$ cocok dengan pengukuran, suku residual
memaksa $u_\theta$ memenuhi persamaan panas dengan $\kappa$ yang sedang ditaksir, dan optimasi menemukan
$\kappa$ yang membuat keduanya konsisten. Solver klasik harus dijalankan berulang kali di dalam lingkar
optimasi untuk masalah yang sama, sedangkan PINN menyelesaikan masalah maju dan invers sekaligus.

Contoh paling terkenal adalah \citet{raissi2020}, yang disinggung di \cref{ch:cv}. Dari video konsentrasi
zat pewarna di sekitar sebuah benda, PINN memulihkan medan kecepatan dan tekanan yang tidak pernah
diukur, karena persamaan transport zat pewarna dan persamaan Navier--Stokes menghubungkan ketiganya.

Contoh yang bisa diperiksa dengan tangan memperlihatkan logika masalah invers. Untuk persamaan panas dengan difusivitas
$\alpha$ yang tidak diketahui di selang $[0,1]$ dengan ujung bernilai nol, fungsi $u(x,t)=e^{-\alpha\pi^2t}\sin(\pi x)$
memenuhi persamaan dan syarat batasnya. Misalkan sensor di tengah batang mencatat $u=1$ pada $t=0$ dan $u=0{,}37$ pada $t=1$.
Residual persamaan hanya bisa nol kalau amplitudo meluruh seperti $e^{-\alpha\pi^2t}$, dan suku data menuntut $e^{-\alpha\pi^2}=0{,}37$,
sehingga $\alpha=-\ln0{,}37/\pi^2\approx0{,}10$. PINN menemukan angka yang sama dengan menurunkan kedua suku loss bersamaan,
tanpa pernah diberi rumus solusinya. Pada masalah nyata, solusinya tidak punya bentuk tertutup, tetapi logikanya sama:
persamaan memberi bentuk, data memberi angka.

Contoh ini juga memperlihatkan bahaya masalah invers. Kalau sensor diletakkan tepat di $x=0$, di mana $\sin(\pi x)=0$,
pengukurannya selalu nol, dan $\alpha$ tidak bisa ditentukan sama sekali. Pengukuran yang tidak peka terhadap parameter
membuat masalahnya \istilah{ill-posed}, dan tidak ada jaringan saraf yang bisa mengatasinya. Memilih letak sensor sama
pentingnya dengan memilih arsitektur.

\section{Syarat batas yang ditegakkan}

Daripada menghukum pelanggaran syarat batas di loss, syarat itu bisa dibangun ke dalam bentuk solusi, seperti $1+tN(t)$ di
contoh tangan tadi. Untuk syarat Dirichlet $u(0)=a$ dan $u(1)=b$ di selang satuan, bentuk
\begin{equation}
u_\theta(x)=a(1-x)+bx+x(1-x)\,N_\theta(x)
\end{equation}
memenuhi kedua syarat untuk jaringan $N_\theta$ apa pun: dua suku pertama adalah garis yang menghubungkan kedua nilai batas,
dan suku ketiga bernilai nol di kedua ujung. Loss tinggal suku residual saja, dan masalah keseimbangan bobot antara residual
dan batas, yang dibahas di bawah, hilang dengan sendirinya. Untuk geometri yang rumit, fungsi $x(1-x)$ diganti fungsi jarak
yang bernilai nol di batas, misalnya signed distance field dari \cref{ch:mlat}. Syarat Neumann, yang menetapkan turunan di
batas, lebih sulit ditegakkan dengan cara ini dan biasanya tetap ditaruh di loss.

Titik collocation juga tidak harus tetap. Pemilihan adaptif berbasis residual menambahkan titik baru di tempat residualnya
paling besar, misalnya dengan menarik sampel dengan peluang sebanding dengan kuadrat residual. Daerah dengan gradien tajam,
seperti lapisan batas atau muka gelombang, mendapat lebih banyak titik, mirip penghalusan mesh adaptif pada solver klasik.

\section{Mengapa PINN sulit dilatih}

Kalau PINN begitu sederhana, mengapa tidak menggantikan semua solver? Pengalaman komunitas selama
beberapa tahun memperlihatkan sejumlah kesulitan yang serius \citep{cuomo2022}.

Kesulitan pertama adalah keseimbangan suku-suku loss. \citet{wang2021} memperlihatkan bahwa gradien suku
residual bisa jauh lebih besar dari gradien suku batas, sehingga jaringan memenuhi persamaan di dalam
domain tetapi mengabaikan syarat batas. Ini masalah condition number dari \cref{ch:optimasi} dalam bentuk
yang khas. Mereka mengusulkan bobot $\lambda$ yang disesuaikan otomatis berdasarkan besar gradien
masing-masing suku. Pendekatan lain memberi setiap titik collocation bobotnya sendiri yang ikut dilatih
dengan arah berlawanan, sehingga titik yang sulit mendapat perhatian lebih \citep{mcclenny2023}.

Skala masalahnya mudah diperkirakan. Misalkan solusinya mengandung komponen $\sin(n\pi x)$. Turunan
keduanya $-(n\pi)^2\sin(n\pi x)$, sehingga galat kecil pada komponen ini diperbesar $(n\pi)^2$ kali dalam residual
dan $(n\pi)^4$ kali dalam kuadratnya. Untuk $n=5$, faktor itu sekitar $6\times10^4$. Suku syarat batas, yang tidak
melibatkan turunan, sama sekali tidak diperbesar. Optimizer melihat lanskap yang hampir seluruhnya ditentukan
residual, dan syarat batas tenggelam.

Kesulitan kedua berasal dari spectral bias (\cref{ch:mlat}). Analisis neural tangent kernel untuk PINN
\citep{wang2022} menunjukkan bahwa komponen solusi dengan frekuensi tinggi dipelajari jauh lebih lambat.
Solusi dengan gradien tajam atau osilasi cepat, yang justru menarik dalam fluida, menjadi sulit. Fourier
features dan aktivasi sinus membantu.

Kesulitan ketiga diperlihatkan \citet{krishnapriyan2021}. Untuk persamaan konveksi dengan koefisien
konveksi yang besar, PINN standar gagal total meski arsitekturnya cukup ekspresif. Masalahnya terletak
pada lanskap loss yang sulit sekali dioptimasi, bukan pada kapasitas jaringan. Melatih bertahap, mulai
dari koefisien kecil lalu dinaikkan, atau memecah domain waktu menjadi potongan yang dilatih berurutan,
memperbaiki keadaan. Gagasan terakhir ini juga menyentuh masalah kausalitas: PINN melatih seluruh rentang
waktu sekaligus, sehingga bisa saja ``memenuhi'' persamaan di waktu akhir sebelum benar di waktu awal,
padahal informasi fisik mengalir dari masa lalu ke masa depan. \citet{wang2024causal} mengusulkan bobot yang
menegakkan urutan itu secara eksplisit. Rentang waktu dibagi menjadi potongan-potongan, dan residual potongan
ke-$i$ diberi bobot $w_i=\exp\big(-\epsilon\sum_{k<i}\Loss_k\big)$. Selama potongan-potongan awal masih punya
residual besar, bobot potongan yang lebih akhir hampir nol, sehingga jaringan dipaksa membereskan masa lalu
lebih dulu.

Ada beberapa trik lain yang sering dipakai: syarat batas yang dibangun langsung ke dalam bentuk solusi,
seperti $1+tN(t)$ pada contoh di atas, sehingga suku batas tidak lagi dibutuhkan; pemilihan titik
collocation secara adaptif, dengan menambah titik di tempat residualnya besar; dan resep optimasi Adam
lalu L-BFGS dari \cref{ch:optimasi}.

\section{Varian}

Beberapa varian mengubah cara persamaan masuk ke dalam loss. Untuk persamaan yang berasal dari prinsip variasi,
residual bisa diganti energi. Persamaan Poisson $-u''=f$ dengan $u(0)=u(1)=0$, misalnya, adalah syarat minimum dari
\begin{equation}
E[u]=\int_0^1\Big(\tfrac12u'(x)^2-f(x)\,u(x)\Big)\dd x .
\end{equation}
Untuk melihatnya, ganggu $u$ menjadi $u+\epsilon v$ dengan $v$ nol di kedua ujung. Turunan $E$ terhadap $\epsilon$ di
$\epsilon=0$ adalah $\int(u'v'-fv)\dd x$. Integrasi parsial memindahkan turunan dari $v$ ke $u$, dan suku batasnya
hilang karena $v$ nol di ujung, sehingga hasilnya $\int(-u''-f)v\dd x$. Supaya nol untuk setiap $v$, haruslah
$-u''=f$. Deep Ritz \citep{e2018} meminimalkan $E[u_\theta]$ secara langsung dengan integrasi Monte Carlo. Keuntungannya,
loss hanya butuh turunan pertama, bukan turunan kedua, sehingga lebih murah dan lebih stabil. Gagasan yang sama
mendasari metode elemen hingga, yang juga bekerja dengan bentuk lemah persamaan.

Untuk domain besar atau solusi yang berubah watak dari satu daerah ke daerah lain, domain bisa dipecah menjadi
beberapa subdomain, masing-masing dengan jaringannya sendiri, yang dijahit dengan suku loss tambahan di
antarmukanya. XPINN \citep{jagtap2020} adalah contoh pendekatan ini, dan ia juga memungkinkan pelatihan paralel.
Gagasannya mirip dekomposisi domain pada solver klasik: masalah besar yang sulit dipecah menjadi masalah kecil
yang lebih mudah.

\section{Penilaian yang jujur}

Untuk masalah maju yang tipikal, misalnya aliran di sekitar benda pada bilangan Reynolds menengah, PINN
umumnya lebih lambat dan kurang akurat daripada solver klasik yang sudah matang. Setiap kasus baru juga
butuh pelatihan baru, karena PINN mempelajari satu solusi, bukan cara menyelesaikan. PINN paling
bernilai di tempat lain: masalah invers dan asimilasi data, masalah dengan data sebaran yang tidak
lengkap, geometri yang sulit di-mesh, dan masalah berdimensi tinggi di mana grid tidak mungkin dibangun
\citep{karniadakis2021}. Untuk mempelajari cara menyelesaikan sebuah keluarga persamaan sekaligus,
kita butuh pendekatan lain: neural operator dan simulator yang dipelajari, tema bab berikutnya.

\section{Perkembangan riset}

Riset PINN beberapa tahun terakhir banyak berfokus pada mengapa pelatihannya sulit dan bagaimana memperbaikinya.
\citet{rathore2024} menghubungkan kesulitan itu dengan condition number: operator diferensial di suku residual membuat
lanskap loss buruk kondisinya, persis contoh skala $(n\pi)^4$ di atas. Mereka membandingkan Adam, L-BFGS, dan gabungan
keduanya, mendapati bahwa gabungan Adam lalu L-BFGS paling baik, dan mengusulkan optimizer orde dua baru yang memperbaiki
kinerja lebih jauh. PirateNets \citep{wang2024pirate} menyerang masalah lain: PINN dengan jaringan yang lebih dalam sering
justru lebih buruk, karena inisialisasi MLP biasa membuat turunan jaringan sulit dilatih. Koneksi residual adaptif mereka
membuat jaringan mulai sebagai jaringan dangkal dan perlahan menjadi dalam selama pelatihan, sehingga kedalaman tambahan
akhirnya memberi akurasi tambahan.

\sumberbab{Bab ini dirangkai dari kuliah Machine Learning: Advanced Techniques (neural field dan loss
yang diinformasikan fisika), kuliah deep learning dasar (autodiff), dan bacaan selama tugas akhir. Bacaan
utama: \citet{raissi2019}, ulasan \citet{karniadakis2021} dan \citet{cuomo2022}, serta studi tentang
kegagalan pelatihan oleh \citet{wang2021}, \citet{wang2022}, dan \citet{krishnapriyan2021}.}
